\scanpage{131}
\diagram{\begin{tikzcd}[ampersand replacement=\&,row sep=small]
X \arrow[d,"f"]\\ S
\end{tikzcd}}
muni d'une section $\sigma$, $f$ propre plat, $f_*(\mathcal O_X)=\mathcal O_S$ \uline{\unsure{universellt|universellement}}.

$G$ schéma en groupes \deleted{\illegible{1}} fini et \deleted{séparable} \inserted{étale} sur $S$
[le cas important est celui où $G=S\times\Gamma$].

$F_G(T)=$ ens des classes à isom près des
fibrés principaux homogènes sur $X_T$, de groupe
$G_T\times_T X_T=(G\times_{S} X_T\deleted{\unsure{\times X|X}}$, \uline{munis} d'une section
marquée le long de $\sigma_T$, i.e. \inserted{d'un isom} $\sigma_T^*(X')\simeq G_T$.

N.B. Tout automorphisme d'un tel torseur
est trivial, \unsure{grâce au|grâce au seul} fait que $G$
est \deleted{\unsure{sép}} étale sur $S$. Il en serait de
$=$ pour $G$ commutatif. Cas général??

Le foncteur $F_G(T)$ en $T$ est-il
représentable ?? Si oui, le \unsure{Sch|$H$} sur $S$
qui le représente est \uline{étale} sur $S$, et
\uline{quasi fini sur} $S$ [car à fibres
\uline{finies}].

% margin note written upside down; the two long strokes cross the note and diagram
\margin{\deleted{4.9, \unsure{4.12|4.10|4.11}.) \quad 4.9.}}
\diagram{\begin{tikzcd}[ampersand replacement=\&,column sep=small,row sep=small]
\& \deleted{C}\arrow[r]\arrow[d,dash] \& \deleted{\overset{\cdot}{C'}}\arrow[r,dash,shift left=1]\arrow[r,dash,shift right=1]\arrow[d] \& \deleted{C''}\arrow[r,dash,shift left=2]\arrow[r,dash]\arrow[r,dash,shift right=2]\arrow[d,dash] \& \deleted{C'''}\arrow[d,dash]\\
\deleted{\text{\uuline{\unsure{épuise|épim.}}}} \& \deleted{\underset{\cdot}{C_0}}\arrow[r,dash] \& \deleted{\underset{\cdot}{C'_0}}\arrow[r,dash,shift left=1]\arrow[r,dash,shift right=1] \& \deleted{C''_0}\arrow[r,dash,shift left=2]\arrow[r,dash]\arrow[r,dash,shift right=2] \& \deleted{C'''_0}\\
\&\& \deleted{\underline{\text{pl. fid}}} \& \deleted{\underline{\text{plat}}} \&\\
\&\& \& \deleted{\underbrace{\hspace{6em}}_{\text{\unsure{exact|exacts}}}} \&
\end{tikzcd}}

Noter que le foncteur $F_G$
est \deleted{\unsure{compati|compatible}} exact à \deleted{\unsure{g}}gauche
pour les morphismes fid. plats
et quasi-compacts.

\deleted{\illegible{2}}
\[
\deleted{G\quad H^1(X,G\quad G\quad G\ \underline{\unsure{\text{tordu}}}}
\]
% archive sheet number: 130
